% prova.tex — conteúdo da prova
%
% A aparência está em provatex.sty e os dados da turma em config.tex.
% Aqui ficam só as orientações e as questões.
%
% Versão do estudante:
%   pdflatex prova.tex
% Versão do professor, com respostas:
%   pdflatex -jobname=prova-respostas "\def\comrespostas{}% prova.tex — conteúdo da prova
%
% A aparência está em provatex.sty e os dados da turma em config.tex.
% Aqui ficam só as orientações e as questões.
%
% Versão do estudante:
%   pdflatex prova.tex
% Versão do professor, com respostas:
%   pdflatex -jobname=prova-respostas "\def\comrespostas{}% prova.tex — conteúdo da prova
%
% A aparência está em provatex.sty e os dados da turma em config.tex.
% Aqui ficam só as orientações e as questões.
%
% Versão do estudante:
%   pdflatex prova.tex
% Versão do professor, com respostas:
%   pdflatex -jobname=prova-respostas "\def\comrespostas{}% prova.tex — conteúdo da prova
%
% A aparência está em provatex.sty e os dados da turma em config.tex.
% Aqui ficam só as orientações e as questões.
%
% Versão do estudante:
%   pdflatex prova.tex
% Versão do professor, com respostas:
%   pdflatex -jobname=prova-respostas "\def\comrespostas{}\input{prova}"
% Ou, para gerar as duas de uma vez:
%   make
\documentclass[brazil, 11pt]{exam}
\usepackage{provatex}
\input{config}

\begin{document}

\maketitle

\begin{multicols}{2}

\begin{orientacoes}
    %\item A prova poderá ser feita em duplas, sem consulta e sem interação entre diferentes duplas.
    \item A prova é individual, sem consulta. Qualquer violação habilitará a invalidação da mesma.
    \item Cada estudante deve entregar sua prova identificada com nome e data.
    \item A prova será recebida somente após 1 hora do seu início.
    %\item As questões 01, 02, 05, 07, 08 e 09 devem estar acompanhadas dos cálculos.
    \item Só serão válidas as respostas fornecidas no cartão-resposta logo abaixo.
    \item É necessário fornecer os 3 dígitos finais da matrícula para devida identificação.
    \item Use caneta azul ou preta para preencher o cartão-resposta.
    \item Deve-se preencher por inteiro o espaço dentro do círculo do cartão-resposta para que a questão seja válida.
    \item Não faça marcações dentro da área do cartão-resposta se não for referente às respostas ou final do número da matrícula.
\end{orientacoes}

\begin{questions}


\begin{question}[1]
Esta é uma questão de exemplo. Normalmente há um texto que pode ser acompanhado de uma figura, uma tabela, uma lista numerada, ou qualquer outra forma de apresentação da informação. \\
Em seguida, são apresentadas as opções para escolha da resposta.

\begin{choices}
    \CorrectChoice Este é o item correto.
    \choice Um item qualquer para ilustrar uma lista de opções.
    \choice Vamos fazer questões objetivas de 5 itens.
    \choice É possível ter uma quantidade maior de itens.
    \choice Um item para encerrar o exemplo.
\end{choices}
\end{question}


\begin{question}[1] % Enem 2012 PPL
\begin{minipage}{\linewidth}\vspace{1ex}\hspace{1em}
Esta é uma questão de exemplo. Normalmente há um texto que pode ser acompanhado de uma figura, uma tabela, uma lista numerada, ou qualquer outra forma de apresentação da informação. \\
Em seguida, são apresentadas as opções para escolha da resposta.

\vspace{1ex}
\begin{choices}
    \CorrectChoice Este é o item correto.
    \choice Um item qualquer para ilustrar uma lista de opções.
    \choice Vamos fazer questões objetivas de 5 itens.
    \choice É possível ter uma quantidade maior de itens.
    \choice Um item para encerrar o exemplo.
\end{choices}

\end{minipage}
\end{question}

\espacocalculos

\listaderespostas

\end{questions}
\end{multicols}
\end{document}
"
% Ou, para gerar as duas de uma vez:
%   make
\documentclass[brazil, 11pt]{exam}
\usepackage{provatex}
% config.tex — dados da turma e da prova
%
% Para uma nova prova, copie este arquivo junto com prova.tex e ajuste
% os dados abaixo. Eles aparecem no cabeçalho e no quadro de orientações.

\curso{Técnico Integrado em Comércio}
\disciplina{Física Complementar}
\atividade{Prova}
\professor{Prof. Paulo Santiago}
\periodo{2025-1}

% Arquivos de imagem
\logo{logo-ifce-bat-11pt.pdf}   % logo do campus, no cabeçalho
\cartaoresposta{gabarito.pdf}   % cartão-resposta gerado pelo EvalBee


\begin{document}

\maketitle

\begin{multicols}{2}

\begin{orientacoes}
    %\item A prova poderá ser feita em duplas, sem consulta e sem interação entre diferentes duplas.
    \item A prova é individual, sem consulta. Qualquer violação habilitará a invalidação da mesma.
    \item Cada estudante deve entregar sua prova identificada com nome e data.
    \item A prova será recebida somente após 1 hora do seu início.
    %\item As questões 01, 02, 05, 07, 08 e 09 devem estar acompanhadas dos cálculos.
    \item Só serão válidas as respostas fornecidas no cartão-resposta logo abaixo.
    \item É necessário fornecer os 3 dígitos finais da matrícula para devida identificação.
    \item Use caneta azul ou preta para preencher o cartão-resposta.
    \item Deve-se preencher por inteiro o espaço dentro do círculo do cartão-resposta para que a questão seja válida.
    \item Não faça marcações dentro da área do cartão-resposta se não for referente às respostas ou final do número da matrícula.
\end{orientacoes}

\begin{questions}


\begin{question}[1]
Esta é uma questão de exemplo. Normalmente há um texto que pode ser acompanhado de uma figura, uma tabela, uma lista numerada, ou qualquer outra forma de apresentação da informação. \\
Em seguida, são apresentadas as opções para escolha da resposta.

\begin{choices}
    \CorrectChoice Este é o item correto.
    \choice Um item qualquer para ilustrar uma lista de opções.
    \choice Vamos fazer questões objetivas de 5 itens.
    \choice É possível ter uma quantidade maior de itens.
    \choice Um item para encerrar o exemplo.
\end{choices}
\end{question}


\begin{question}[1] % Enem 2012 PPL
\begin{minipage}{\linewidth}\vspace{1ex}\hspace{1em}
Esta é uma questão de exemplo. Normalmente há um texto que pode ser acompanhado de uma figura, uma tabela, uma lista numerada, ou qualquer outra forma de apresentação da informação. \\
Em seguida, são apresentadas as opções para escolha da resposta.

\vspace{1ex}
\begin{choices}
    \CorrectChoice Este é o item correto.
    \choice Um item qualquer para ilustrar uma lista de opções.
    \choice Vamos fazer questões objetivas de 5 itens.
    \choice É possível ter uma quantidade maior de itens.
    \choice Um item para encerrar o exemplo.
\end{choices}

\end{minipage}
\end{question}

\espacocalculos

\listaderespostas

\end{questions}
\end{multicols}
\end{document}
"
% Ou, para gerar as duas de uma vez:
%   make
\documentclass[brazil, 11pt]{exam}
\usepackage{provatex}
% config.tex — dados da turma e da prova
%
% Para uma nova prova, copie este arquivo junto com prova.tex e ajuste
% os dados abaixo. Eles aparecem no cabeçalho e no quadro de orientações.

\curso{Técnico Integrado em Comércio}
\disciplina{Física Complementar}
\atividade{Prova}
\professor{Prof. Paulo Santiago}
\periodo{2025-1}

% Arquivos de imagem
\logo{logo-ifce-bat-11pt.pdf}   % logo do campus, no cabeçalho
\cartaoresposta{gabarito.pdf}   % cartão-resposta gerado pelo EvalBee


\begin{document}

\maketitle

\begin{multicols}{2}

\begin{orientacoes}
    %\item A prova poderá ser feita em duplas, sem consulta e sem interação entre diferentes duplas.
    \item A prova é individual, sem consulta. Qualquer violação habilitará a invalidação da mesma.
    \item Cada estudante deve entregar sua prova identificada com nome e data.
    \item A prova será recebida somente após 1 hora do seu início.
    %\item As questões 01, 02, 05, 07, 08 e 09 devem estar acompanhadas dos cálculos.
    \item Só serão válidas as respostas fornecidas no cartão-resposta logo abaixo.
    \item É necessário fornecer os 3 dígitos finais da matrícula para devida identificação.
    \item Use caneta azul ou preta para preencher o cartão-resposta.
    \item Deve-se preencher por inteiro o espaço dentro do círculo do cartão-resposta para que a questão seja válida.
    \item Não faça marcações dentro da área do cartão-resposta se não for referente às respostas ou final do número da matrícula.
\end{orientacoes}

\begin{questions}


\begin{question}[1]
Esta é uma questão de exemplo. Normalmente há um texto que pode ser acompanhado de uma figura, uma tabela, uma lista numerada, ou qualquer outra forma de apresentação da informação. \\
Em seguida, são apresentadas as opções para escolha da resposta.

\begin{choices}
    \CorrectChoice Este é o item correto.
    \choice Um item qualquer para ilustrar uma lista de opções.
    \choice Vamos fazer questões objetivas de 5 itens.
    \choice É possível ter uma quantidade maior de itens.
    \choice Um item para encerrar o exemplo.
\end{choices}
\end{question}


\begin{question}[1] % Enem 2012 PPL
\begin{minipage}{\linewidth}\vspace{1ex}\hspace{1em}
Esta é uma questão de exemplo. Normalmente há um texto que pode ser acompanhado de uma figura, uma tabela, uma lista numerada, ou qualquer outra forma de apresentação da informação. \\
Em seguida, são apresentadas as opções para escolha da resposta.

\vspace{1ex}
\begin{choices}
    \CorrectChoice Este é o item correto.
    \choice Um item qualquer para ilustrar uma lista de opções.
    \choice Vamos fazer questões objetivas de 5 itens.
    \choice É possível ter uma quantidade maior de itens.
    \choice Um item para encerrar o exemplo.
\end{choices}

\end{minipage}
\end{question}

\espacocalculos

\listaderespostas

\end{questions}
\end{multicols}
\end{document}
"
% Ou, para gerar as duas de uma vez:
%   make
\documentclass[brazil, 11pt]{exam}
\usepackage{provatex}
% config.tex — dados da turma e da prova
%
% Para uma nova prova, copie este arquivo junto com prova.tex e ajuste
% os dados abaixo. Eles aparecem no cabeçalho e no quadro de orientações.

\curso{Técnico Integrado em Comércio}
\disciplina{Física Complementar}
\atividade{Prova}
\professor{Prof. Paulo Santiago}
\periodo{2025-1}

% Arquivos de imagem
\logo{logo-ifce-bat-11pt.pdf}   % logo do campus, no cabeçalho
\cartaoresposta{gabarito.pdf}   % cartão-resposta gerado pelo EvalBee


\begin{document}

\maketitle

\begin{multicols}{2}

\begin{orientacoes}
    %\item A prova poderá ser feita em duplas, sem consulta e sem interação entre diferentes duplas.
    \item A prova é individual, sem consulta. Qualquer violação habilitará a invalidação da mesma.
    \item Cada estudante deve entregar sua prova identificada com nome e data.
    \item A prova será recebida somente após 1 hora do seu início.
    %\item As questões 01, 02, 05, 07, 08 e 09 devem estar acompanhadas dos cálculos.
    \item Só serão válidas as respostas fornecidas no cartão-resposta logo abaixo.
    \item É necessário fornecer os 3 dígitos finais da matrícula para devida identificação.
    \item Use caneta azul ou preta para preencher o cartão-resposta.
    \item Deve-se preencher por inteiro o espaço dentro do círculo do cartão-resposta para que a questão seja válida.
    \item Não faça marcações dentro da área do cartão-resposta se não for referente às respostas ou final do número da matrícula.
\end{orientacoes}

\begin{questions}


\begin{question}[1]
Esta é uma questão de exemplo. Normalmente há um texto que pode ser acompanhado de uma figura, uma tabela, uma lista numerada, ou qualquer outra forma de apresentação da informação. \\
Em seguida, são apresentadas as opções para escolha da resposta.

\begin{choices}
    \CorrectChoice Este é o item correto.
    \choice Um item qualquer para ilustrar uma lista de opções.
    \choice Vamos fazer questões objetivas de 5 itens.
    \choice É possível ter uma quantidade maior de itens.
    \choice Um item para encerrar o exemplo.
\end{choices}
\end{question}


\begin{question}[1] % Enem 2012 PPL
\begin{minipage}{\linewidth}\vspace{1ex}\hspace{1em}
Esta é uma questão de exemplo. Normalmente há um texto que pode ser acompanhado de uma figura, uma tabela, uma lista numerada, ou qualquer outra forma de apresentação da informação. \\
Em seguida, são apresentadas as opções para escolha da resposta.

\vspace{1ex}
\begin{choices}
    \CorrectChoice Este é o item correto.
    \choice Um item qualquer para ilustrar uma lista de opções.
    \choice Vamos fazer questões objetivas de 5 itens.
    \choice É possível ter uma quantidade maior de itens.
    \choice Um item para encerrar o exemplo.
\end{choices}

\end{minipage}
\end{question}

\espacocalculos

\listaderespostas

\end{questions}
\end{multicols}
\end{document}
